\subsection{History and the ingredients of the proof}

The random-cluster representation of Fortuin and Kasteleyn connects
percolation and Potts-type spin systems through a weight on edge
subsets~\cite{FortuinKasteleyn1972}. On a random planar map, the map
itself is also sampled, so the cluster weight changes the underlying
geometry. Sheffield's inventory model encodes this joint randomness by
two types of burgers, two fixed order types, and flexible orders that
take the freshest available burger~\cite{Sheffield2016}. Its parameter
$p=\sqrt q/(\sqrt q+2)$ places $q>4$ exactly in the regime $p>1/2$.
Sheffield's unconditioned walk theorem identifies the transition at
$q=4$: above it, the limiting inventory fluctuations are
one-dimensional~\cite[Theorem~2.5]{Sheffield2016}.

The finite-volume CRT prediction appears in Sheffield's
appendix~\cite{Sheffield2016}. Feng restates it as
Conjecture~1.2 and proves an infinite-volume
counterpart~\cite{Feng2026}: his Theorem~1.3 concerns an infinite FK map
and convergence to an infinite continuum random tree in the
\emph{local} Gromov--Hausdorff--Prokhorov topology.
Theorem~\ref{thm:main} realizes the finite-volume prediction in the
compact topology, with the degree probability measure and all integer
edge sizes. It is proved directly under the finite law, not deduced
from the infinite-map limit.

The contrast with the other FK regimes is geometric. Companion results
for $0<q<4$ give Liouville quantum gravity sphere limits in the
metric-measure and conformal settings~\cite{SubcriticalMetric2026,SubcriticalConformal2026};
the critical companion treats the sphere and FK-map geometry at
$q=4$~\cite{CriticalFK2026}. Above four the limit here is instead a
tree. These companion results provide context and are not inputs to
the proof.

The finite-map encoding has a longer combinatorial ancestry. Mullin
enumerated tree-rooted maps~\cite{Mullin1967}, and Bernardi gave an
explicit bijective account using shuffles of two parenthesis
systems~\cite{Bernardi2007}. Bernardi's embedding activities and
subgraph correspondences relate these encodings to the Tutte
polynomial~\cite{Bernardi2008,Bernardi2008Subgraphs}; Sheffield combines
this structure with flexible orders to obtain the FK
weights~\cite[Section~4]{Sheffield2016}. We rederive the needed finite
identity with our dart-rooting and activity conventions in
Section~\ref{word:section}. The derivation fixes the exact conditioned
map law, including loops, links, and the trivial map used in the
generating function.

The inventory estimates below also build on specific earlier work.
Sheffield analyzes the first backward surviving burger and computes
the associated mean reduced length in
\cite[Section~3.1, Lemma~3.1]{Sheffield2016}. We use the same stationary
covariance and stopping method to prove the strict drift estimate
needed here. The reduced ladder excursions, balanced burger frequencies,
positive flexible-order frequency, and enclosing empty intervals of
\cite[Section~3.3]{Sheffield2016} are close predecessors of our cut
construction. We use these ideas to obtain the uniform matching
probability and generating-function lower bound needed for the
finite-map argument; both estimates are proved in
Section~\ref{ladder:section}.

The block decomposition itself goes back to
Tutte~\cite[Section~6]{Tutte1963}. Addario-Berry makes its exact
corner-indexed tree, including empty insertions and the $2n+1$ node
count, explicit for uniform maps~\cite[Section~2 and Proposition~3.1]{AddarioBerry2019}.
Stufler treats block-weighted maps in the general framework of enriched
trees and describes their metrics in terms of
corners~\cite[preprint Sections~6.1.5 and~6.8.3]{Stufler2020}.
Exponential tilting preserves a simply generated tree law conditional
on its size; Janson gives the associated offspring moment
formulas~\cite[Section~4]{Janson2012}.
The CRT limits for subcritical graph classes of Panagiotou, Stufler,
and Weller~\cite{PanagiotouStuflerWeller2016}, and the block-map results
in~\cite[preprint Theorems~6.60 and~6.62]{Stufler2020}, establish the principle
that graph distance is asymptotic to a constant times ancestral depth.
Their analytic hypotheses supply exponential moments. Our endpoint
argument establishes only finite offspring variance, so we retain a
path estimate that requires no exponential or third moment.

A recent result of Stufler~\cite[Theorem~1.2 and Lemma~9.1]{Stufler2026}
proves finite-variance degree-measure Gromov--Hausdorff--Prokhorov CRT
limits for maps whose block weights are functions of the block's edge
count. It also uses the exact corner-to-degree correspondence. In that
model blocks of a given size are uniform, and the metric argument uses
diameter estimates for uniform nonseparable maps. The FK weight of a
block instead depends on its shape, not only on its edge count. We must
prove finite variance for these weights and control their particular
metric marks.
The quadratic inverse criterion in Section~\ref{block:section} and the
finite-second-moment path argument in Section~\ref{tree:section}
provide those steps.

Finally, the limiting tree and the probabilistic tools have classical
origins. Aldous introduced the Brownian continuum random tree and its
finite-variance branching-process limit~\cite{Aldous1993}; Duquesne
developed the stable-domain extension~\cite{Duquesne2003}.
Marckert and Mokkadem proved joint depth-first process convergence to
the same excursion under stronger moment
hypotheses~\cite{MarckertMokkadem2003}. We use Broutin and Marckert's
joint coding theorem and conditioned-degree lemma in a formulation
that gives the height, contour, and walk limits jointly under a finite
second moment, through every attainable size~\cite{BroutinMarckert2014}.
The attainable tree sizes here are exactly the odd integers. The marked
ancestral-path calculation uses the size-biased spine method of
Lyons, Pemantle, and Peres~\cite[Section~2]{LyonsPemantlePeres1995}.
The truncation that makes it uniform over all nodes is proved below.
